% This is the Foreword to the Hodgkin-Huxley neuron model V2.0 book

\chapter*{Foreword\label{ch:Foreword}}
\addcontentsline{toc}{chapter}{Foreword} 


Hodgkin and Huxley's classic (V1.0) "model", which fundamentally changed our ideas about the physiology of neurons and which has served as the basis for describing the functioning of neurons for several decades, can hardly be overestimated; it is truly the crown jewel of neuroscience. However, its authors, quite realistically, assessed that the physical picture behind the observations they mathematically formulated could be partially incorrect. Already in their original communication, they drew attention to the fact that their equations did not fully describe their observations. Later they added that there exists observed evidence
that does not fit into the picture.
Due to the lack of certain knowledge, they were forced to formulate ad hoc assumptions and they admitted that the formalism of the equations could describe more than one physical process.
Some of the assumptions proved to be a brilliant vision, but in quite a few cases it turned out that they did not always put the right physical background behind their observations
which were generally correctly described mathematically, and that their model designed for mechanical computers had obvious limitations.
Already at the time of the introduction of the model, doubts arose as to whether the purely electrical model could fully describe the observations known up to that time, and later, as the accuracy, directionality and detail of the observations increased, it became necessary to introduce more and more ad hoc assumptions in order to maintain the consistency of the model assuming purely electrical operation.


% \vbox{
%"\textit{Though this be madness, yet there is method in't.}"
%
%\hfill Shakespeare: Hamlet: Act 2 Scene 2
%	\vskip 2 in
%   \null
%    \centering
%    {\LARGE \@title \par}%
%%    \vskip 1.0em%
%%    {\large \@subtitle \par}%
%    \vskip 10em%
%    {\large \@author \par}%
%    \vskip 1.5em%
%    {\large \@date \par}%       % Set date in \large size.
%    \vskip 3em%
%   \null
%   }%

%Az elmúlt évtizedekben számos (V1.0x) modell jött létre, több-kevesebb sikerrel módosítva és kiegészítve ezt a tisztán elektromos, diszciplináris, őseredeti modellt (ami végső soron \gls{HH} követőinek 
%tevékenysége nyomán vált modellé). Születtek olyan (V1.y) modellek is, amelyek 
%teljesen más diszciplina keretében próbálják a neuront modellezni,
%főként annak termodinamikai, mechanikai, optikai, stb. jelenségeit értelmezni.
%Mindezen modellek közös sajátsága, hogy nem tudnak teljes körű
%(valamennyi megfigyelést ellentmondásmentesen leíró) modellt nyújtani.
%Vannak olyan modellek is, amelyek nem is törekszenek arra, hogy
%egy fizikai működési mechanizmust próbáljanak meg leírni;
%megelégszenek egy olyan matematikai formalismussal, amelyik
%a tapasztalattal megegyező viselkedést mutat; anélkül, hogy
%a 'miért', 'milyen mechanizmussal' kérdésekre akár csak megpróbálnának
%választ adni. 

In the past decades, numerous (V1.0x) models have been created, modifying and supplementing this purely electrical, disciplinary, original model with more or less success (which ultimately became a model due to the activities of \gls{HH}'s followers).
There have also been (V1.y) models that try to model the neuron within the framework of an entirely different discipline, mainly interpreting its thermodynamic, mechanical, optical, etc. phenomena.
The common feature of all these models is that they cannot provide a complete model (describing all observations without contradiction).
There are also models that do not even strive to describe a physical operating mechanism; they are satisfied with providing a mathematical formalism that shows behavior that is consistent with experience;
without even trying to answer the questions 'why' and 'by what mechanism'.

Due to the lack of the needed expertize in physics,
biophysics introduced the idea of the "whatnot"~\cite{Schrodinger:1992} protein mechanisms 
for the unexplainable forces that totally misguided physiology and neuroscience.
Biology has no laws of motion in the sense as physics has Newton's laws.
The many speculative hypotheses hyde that science can describe
living matter, but needs a different approach.
Claims such as "pumps [by using protein mechanisms] \dots transport ions \textit{against their
electrical and chemical gradients}"~\cite{PrinciplesNeuralScience:2013}, page 101, 
or "the flux of ions through ion channels is passive,
requiring \textit{no expenditure of metabolic energy} by the
channels."~\cite{PrinciplesNeuralScience:2013}, page 107, are as scientific as in the pre-Newtonian age,
there was a notion related to Aristotelian philosophy. This concept gave credence to the “theory” that the celestial spheres and, later on, the planets were pushed or moved by celestial movers, celestial intelligences, or angels against physical forces, without continuously investing energy. Inventing the universal 
gravitational law made the complex interplay of those spheres obsolete,
although, for the first look, it seemed to be unrelated.

Describing life science is still in the pre-Newtonian age: it has a lot of "spheres"
described by complicated theories that contradict each other,
and does not have its laws of motion, in the sense as Newton's Laws in physics. What is worse, it did not invent and test the laws describing
living matter. Instead, it borrowed the disciplinary laws of physics,
without recognizing that its different discipline has different
conditions and requires different approximations; that is similar, but differently used laws. 
\index{law of motion}
Without questioning the merits of the V1.0 model, the experimental experience accumulated in recent decades, the dizzying development of technology and instrumental measurements, and the ever-increasing demand for the results of neuroscience have made it necessary to develop the (V2.0) model. The model is necessarily cross-disciplinary; no single physical discipline alone allows for a complete and consistent interpretation of the phenomena of the living organism.

%Nem vonva kétségbe a V1.0 modell érdemeit, az elmúlt évtizedekben 
%felhalmozott kísérleti tapasztalatok, a technika és a műszeres mérésések  
%szédítő mértékű fejlődése, az idegtudomány eredményei iránti igény
%mind erősebb növekedése szükségessé teszi (V2.0) modell kidolgozását.
%A modell cross-disciplináris; egyedül egyik fizikai disciplina sem teszi
%lehetővé az élő szervezet jelenségeinek teljes körű és 
%ellentmondásmentes értelmezését.
%







